% ID: SEL-D0064-OB-E01
% ESTADO: cotejada. Segunda pasada visual y borrador compilado revisados.
% Fuente: Selectividad/sel_2010_matematicas_aplicadas/mat_apli/mat_apli_sept/mat_apli_sept/EXAMEN-ANDALUCÍA-3.pdf
% SHA-256: 9c3d2d628115a2f20f17b47171106ddd90f214b5f6804bdf9a7dcb7037c75efd
% Página: 2; opción B; ejercicio original 1.

Sean las matrices:
\[
P=\begin{pmatrix}1&2\\a&0\end{pmatrix},\quad Q=\begin{pmatrix}1&1&5\\8&4&b\end{pmatrix}\quad\text{y}\quad R=\begin{pmatrix}c&d&6\\10&10&50\end{pmatrix}.
\]
\begin{enumerate}[label=\alph*)]
\item \textbf{(1 punto)} Calcule, si es posible, $P\cdot Q$ y $Q\cdot P$, razonando la respuesta.
\item \textbf{(1.5 puntos)} \textquestiondown{}Cuánto deben valer las constantes $a$, $b$, $c$ y $d$ para que $P\cdot2Q=R$?
\end{enumerate}
